\section{Supporting Gaussian model and exit-time
results}\label{draft-supporting-gaussian-model-and-exit-time-results}

Part of the \href{DRAFT.md}{manuscript}. This chapter retains the
original Gaussian model, propositions 1--12, tolerance regimes and
supporting discussion. Its original section numbers are retained for the
proof references in \url{gaussian_proofs.md}. The independent prediction
problem and current statistical results are defined in the main
manuscript; these Gaussian mixture results are supporting geometry, not
additional assumptions on the common-risk tilt family.

\subsection{3. Model, information, and risk
safety}\label{draft-model-information-and-risk-safety}

\subsubsection{3.1. Anchored mixture and the deployed
classifier}\label{draft-anchored-mixture-and-the-deployed-classifier}

Let the target feature and label satisfy \[
Z=\delta+Ym+\sigma\xi,\qquad
\Pr(Y=1)=p,\qquad \xi\sim N(0,1),
\] where \(Y\in\{-1,+1\}\), \(0<p<1\), and \(m,\sigma,w>0\). The
positive class orientation and the parameters \(m,\sigma\) are supplied
by a source anchor. The translation \(\delta\) and target prior \(p\)
characterize the shift. An oracle analysis may use them to describe the
dynamics; an unlabeled deployment rule would have to estimate or
otherwise handle their uncertainty.

The frozen classifier has threshold \(b_0=0\). Its prediction is
\(\operatorname{sign}(Z-b)\), and its positive-class confidence is \[
s(w(Z-b)),\qquad s(u)=(1+e^{-u})^{-1}.
\] Only \(b\) adapts. The feature representation and positive slope
\(w\) remain fixed.

Center the feature and threshold: \[
z=Z-\delta,\qquad x=b-\delta,\qquad x_0=-\delta.
\] The centered feature density is \[
f_p(z)=p\varphi_\sigma(z-m)+(1-p)\varphi_\sigma(z+m),
\] where \(\varphi_\sigma(z)=\sigma^{-1}\varphi(z/\sigma)\), and
\(\varphi,\Phi\) are the standard normal density and CDF. Translation
changes the initialization and the frozen-risk comparator, but not the
centered entropy vector field.

\subsubsection{3.2. Actual entropy
updates}\label{draft-actual-entropy-updates}

Let \[
a(u)=s(u)(1-s(u)),\qquad
h(u)=-s(u)\log s(u)-(1-s(u))\log(1-s(u)).
\] Define expected entropy and rightward velocity by \[
F_p(x)=\mathbb E_p h(w(z-x)),\qquad v_p(x)=-F_p'(x).
\] The rightward update from one sample is \[
U(x,z)=-w^2(z-x)a(w(z-x)),\qquad \mathbb E_pU(x,z)=v_p(x).
\] Thus population gradient descent is \[
\bar x_{t+1}=\bar x_t+\eta v_p(\bar x_t),
\] and fresh-batch SGD is \[
x_{t+1}=x_t+\frac{\eta}{n}\sum_{i=1}^n U(x_t,z_{t,i}).
\] The batch size is \(n\). Every batch is independent of the past, and
all samples have the same target distribution. Labels are not read by
the update. Dependent reuse of observations, changing priors, momentum,
and feature learning are outside this recursion. The visual experiments
below instead draw independently with replacement from a fixed feature
pool: conditional on that pool, the samples are iid from its empirical
distribution. This preserves sampling independence but does not make
that non-Gaussian empirical distribution satisfy the Gaussian model
assumptions.

\subsubsection{3.3. The comparator and the first unsafe
iterate}\label{draft-the-comparator-and-the-first-unsafe-iterate}

Classification risk on an independent target draw is \[
r_p(x)=p\Phi((x-m)/\sigma)+(1-p)\Phi((-x-m)/\sigma).
\] The corresponding uncentered risk is \(R(b)=r_p(b-\delta)\). For
tolerance \(\epsilon\ge0\), define \[
T_\epsilon=\inf\{t\ge0:r_p(x_t)>r_p(x_0)+\epsilon\},
\qquad \inf\varnothing=\infty.
\] The comparator is the frozen classifier evaluated on this same target
distribution, not its source training risk. The strict inequality
specifies the first unsafe iterate. The risk-safe set is \[
\mathcal S_\epsilon=\{x:r_p(x)\le r_p(x_0)+\epsilon\}.
\]

A deployment rule intended to guarantee safety must reject an unsafe
proposal before it is used. Low average risk over time does not imply
that every deployed iterate is safe. Nor does an exit to the more
frequent class represent safety: either boundary of
\(\mathcal S_\epsilon\) is a violation.

\subsection{4. Exact identities and entropy-kernel
geometry}\label{draft-exact-identities-and-entropy-kernel-geometry}

\subsubsection{Proposition 1. Risk geometry and uniform update
bounds}\label{draft-proposition-1.-risk-geometry-and-uniform-update-bounds}

The risk has a unique finite minimizer \[
x_*(p)=\frac{\sigma^2}{2m}\log\frac{1-p}{p}.
\] It decreases strictly left of \(x_*\), increases strictly right of
\(x_*\), and satisfies \[
r_p(-\infty)=1-p,\qquad r_p(+\infty)=p.
\] Moreover, \[
|U(x,z)|\le M:=w/4,\qquad
|\partial_xU(x,z)|\le L:=w^2/4.
\] If \(0<\eta<2/L\), every nonstationary population update strictly
decreases \(F_p\).

\textbf{Proof.} See \href{gaussian_proofs.md}{Gaussian proofs}, under
the same result number.

\subsubsection{Lemma 2. The entropy kernel is
log-concave}\label{draft-lemma-2.-the-entropy-kernel-is-log-concave}

The function \(h\) is positive, even, integrable, and log-concave. If \[
K(t)=\mathbb E h(w(t+\sigma\xi)),
\] then \(K\) is positive and even, and \(K'(t)<0\) for \(t>0\).
Consequently, \[
F_p(x)=pK(x-m)+(1-p)K(x+m).
\]

\textbf{Proof.} See \href{gaussian_proofs.md}{Gaussian proofs}, under
the same result number.

We use the standard convolution closure property: if two nonnegative
integrable functions are log-concave, their convolution is log-concave
{[}7{]}. The model-specific task is verifying the kernel and the mixture
density, rather than claiming this closure result as a contribution.

\subsection{5. Distribution structure determines population
behavior}\label{draft-distribution-structure-determines-population-behavior}

\subsubsection{Theorem 3. Weak separation gives a unique unstable
entropy
peak}\label{draft-theorem-3.-weak-separation-gives-a-unique-unstable-entropy-peak}

Assume \(0<m\le\sigma\). There is a unique stationary point \(q_p\) of
\(F_p\), and it is a strict global maximum. Its drift signs are \[
v_p(x)<0\quad(x<q_p),\qquad v_p(x)>0\quad(x>q_p).
\] The sign of \(q_p\) is the sign of \(p-1/2\). For every \(\eta>0\),
population iterates initialized right of \(q_p\) increase to
\(+\infty\); iterates initialized left of it decrease to \(-\infty\).
Initialization at \(q_p\) stays fixed.

\textbf{Proof.} See \href{gaussian_proofs.md}{Gaussian proofs}, under
the same result number.

Escape monotonicity holds for every positive step size. Entropy descent
additionally requires a condition such as \(0<\eta<2/L\).

\subsubsection{Corollary 4. An explicit mixed-shift rebound
region}\label{draft-corollary-4.-an-explicit-mixed-shift-rebound-region}

Fix \(\epsilon\ge0\). Suppose \[
0<m\le\sigma,\qquad
\Phi(-m/\sigma)+\epsilon<p<1/2,\qquad -x_*(p)<\delta\le0,
\] and \[
0<\eta<\min\left\{\frac8{w^2},
\frac{4[x_*(p)+\delta]}{w}\right\}.
\] Then the first population update strictly reduces classification
risk, every population update strictly reduces entropy, and
\(T_\epsilon<\infty\). The iterates tend to \(+\infty\) and risk tends
to \(p\). Taking \(m<\sigma,\delta<0\) with strict inequalities gives a
nonempty open set with genuine translation and prior shift. Reflection
\((x,\delta,p)\mapsto(-x,-\delta,1-p)\) gives the other orientation.

\textbf{Proof.} See \href{gaussian_proofs.md}{Gaussian proofs}, under
the same result number.

For \(p<1/2\), weak separation also yields the following classification.

{\def\LTcaptype{table} % do not increment counter
\begin{longtable}[]{@{}
  >{\raggedright\arraybackslash}p{(\linewidth - 4\tabcolsep) * \real{0.3333}}
  >{\raggedright\arraybackslash}p{(\linewidth - 4\tabcolsep) * \real{0.3333}}
  >{\raggedright\arraybackslash}p{(\linewidth - 4\tabcolsep) * \real{0.3333}}@{}}
\toprule\noalign{}
\begin{minipage}[b]{\linewidth}\raggedright
Initialization
\end{minipage} & \begin{minipage}[b]{\linewidth}\raggedright
Threshold motion
\end{minipage} & \begin{minipage}[b]{\linewidth}\raggedright
Risk behavior
\end{minipage} \\
\midrule\noalign{}
\endhead
\bottomrule\noalign{}
\endlastfoot
\(x_0<q_p\) & To \(-\infty\) & Immediately worsens; limit \(1-p\) \\
\(x_0=q_p\) & Fixed & No population adaptation \\
\(q_p<x_0<x_*\) & To \(+\infty\) & Initially improves for a sufficiently
small step; then rises \\
\(x_0\ge x_*\) & To \(+\infty\) & Strictly worsens; limit \(p\) \\
\end{longtable}
}

On a rightward trajectory, \(T_\epsilon<\infty\) exactly when
\(r_p(x_0)+\epsilon<p\). On a leftward trajectory the corresponding
criterion uses \(1-p\). Thus collapse toward one predicted class does
not automatically imply crossing the risk of an already poor
initialization.

\subsubsection{Theorem 5. A sufficiently separated safe population
basin}\label{draft-theorem-5.-a-sufficiently-separated-safe-population-basin}

Set \(p=1/2\), \(a_m=m/\sigma>1\), and \[
C_0=\int h(u)\,du=\pi^2/3,\qquad
C_2=\int u^2h(u)\,du=7\pi^4/45.
\] If \[
(w\sigma)^2>
\frac{3C_2}{2C_0}\frac{e^{a_m^2/2}}{a_m^2-1},
\] then \(F_{1/2}''(0)>0\). There is an \(r>0\) such that, for
\(0<|\delta|<r\) and \(0<\eta\le1/L\), population iterates converge
monotonically to zero and their risks strictly decrease to Bayes risk.

For any fixed nonzero initialization strictly inside this basin and
\(\eta<1/L\), a sufficiently small open perturbation of the parameters
permits \(p\ne1/2,\delta\ne0\): the iterates converge to a finite
entropy minimum \(x_e\), and \[
r_p(x_t)\le r_p(x_0)\quad\text{for every }t,\qquad
r_p(x_e)<r_p(x_0).
\]

\textbf{Proof.} See \href{gaussian_proofs.md}{Gaussian proofs}, under
the same result number.

This is a local sufficient condition, not the exact boundary between all
stable and unstable regimes. It does not guarantee that fixed-step SGD
remains in the basin forever.

\subsection{6. Deterministic safety duration near class
balance}\label{draft-deterministic-safety-duration-near-class-balance}

Assume \(0<m\le\sigma\) and write \(p_\rho=(1-\rho)/2\). Subscripts
\(\rho\) denote the corresponding prior. Define \[
\beta=-K'(m)>0,\qquad \lambda=-K''(m)>0,\qquad
k_0=\sigma^2/m,\qquad A=1+\eta\lambda.
\] Positivity of \(\lambda\) follows from the balanced density \(f_0\):
for \(z>0\), \[
(\log f_0)'(z)=\frac{-z+m\tanh(mz/\sigma^2)}{\sigma^2}<0.
\] Since \(h_w'(z)<0\), integration by parts gives \[
F_0''(0)=-\int h_w'(z)f_0'(z)\,dz<0.
\] Near the origin, \[
v_\rho(x)=\beta\rho+\lambda x+O(x^2+|\rho x|+\rho^2).
\] The implicit function theorem then gives \[
q_\rho=-\frac\beta\lambda\rho+O(\rho^2),\qquad
\lambda_\rho=v_\rho'(q_\rho)=\lambda+O(\rho).
\]

\subsubsection{Theorem 6. Fixed and shrinking tolerances have different
time
scales}\label{draft-theorem-6.-fixed-and-shrinking-tolerances-have-different-time-scales}

Fix \(\eta>0\) and initialize \(x_0=\alpha_0\rho\), where \[
-\beta/\lambda<\alpha_0<k_0.
\] For fixed \[
0<\epsilon<1/2-\Phi(-m/\sigma),
\] the population crossing time satisfies \[
T_\epsilon(\rho)=\frac{\log(1/\rho)}{\log(1+\eta\lambda)}+O(1),
\qquad \rho\downarrow0.
\] If also \[
\eta<\min\{2/L,\;(k_0-\alpha_0)/(\beta+\lambda\alpha_0)\},
\] the first update improves risk and every update reduces entropy.

For \(\epsilon=c\rho^2\), \(c\ge0\), the crossing time is bounded as
\(\rho\downarrow0\). Put \[
B=\frac{m\varphi(m/\sigma)}{\sigma^3},\qquad
u_t=(\alpha_0+\beta/\lambda)A^t-\beta/\lambda,
\] \[
u_c=k_0+\sqrt{(\alpha_0-k_0)^2+2c/B}.
\] If no integer \(t\) satisfies \(u_t=u_c\), then, for sufficiently
small \(\rho>0\), \[
T_{c\rho^2}=
\left\lfloor
\frac{\log[(u_c+\beta/\lambda)/(\alpha_0+\beta/\lambda)]}
{\log A}\right\rfloor+1.
\] At an integer equality \(u_j=u_c\), the first crossing is \(j\) or
\(j+1\), depending on higher-order terms.

\textbf{Proof.} See \href{gaussian_proofs.md}{Gaussian proofs}, under
the same result number.

At exact balance with \(x_0=0\), the population recursion remains fixed
forever. This does not contradict either limiting regime. The available
risk improvement under a small prior perturbation is itself only
\(O(\rho^2)\). A longer fixed-tolerance horizon need not mean a larger
useful improvement.

\subsection{7. Direct first-exit bounds in a positive-drift
corridor}\label{draft-direct-first-exit-bounds-in-a-positive-drift-corridor}

Write SGD as \[
x_{t+1}=x_t+\eta[v_p(x_t)+\xi_t],\qquad
\mathbb E[\xi_t\mid\mathcal F_t]=0,
\] where \(\mathcal F_t\) contains the history before the new batch.
Proposition 1 and Hoeffding's lemma imply \[
\mathbb E[e^{s\xi_t}\mid\mathcal F_t]
\le\exp(s^2M^2/(2n)).
\] The conditional formulation permits state-dependent noise.

\subsubsection{Theorem 7. Two-sided first-exit
control}\label{draft-theorem-7.-two-sided-first-exit-control}

Suppose \(\mathcal S_\epsilon=[a,c]\) is finite, \(a<x_0<c\), and \[
0<\nu\le v_p(x)\le V\qquad(x\in[a,c]).
\] Let \(d=x_0-a\), \(D=c-x_0\), and \(\tau=T_\epsilon\). Then \[
\Pr(x_\tau<a,\tau<\infty)\le
\exp\left(-\frac{2n\nu d}{\eta M^2}\right).
\] For integers \(t\ge1\), \[
\Pr(\tau\le t,x_\tau>c)
\le\exp\left[-\frac{n(D-\eta Vt)^2}{2\eta^2M^2t}\right]
\quad\text{if }D>\eta Vt,
\] and \[
\Pr(\tau>t)
\le\exp\left[-\frac{n(\eta\nu t-D)^2}{2\eta^2M^2t}\right]
\quad\text{if }\eta\nu t>D.
\] Calling the three bounds \(\alpha_L,\alpha_E(t),\alpha_D(t)\),
respectively, \[
\Pr(t_-<T_\epsilon\le t_+,\ x_{T_\epsilon}>c)
\ge1-\alpha_L-\alpha_E(t_-)-\alpha_D(t_+).
\]

\textbf{Proof.} See \href{gaussian_proofs.md}{Gaussian proofs}, under
the same result number.

The theorem allows backward stochastic steps and overshoot. Its price is
a positive drift bound on the entire safe interval. It cannot cover a
critical safe interval containing the entropy peak.

For comparison, a standard global tracking argument gives, for fixed
\(T\ge1\) and failure level \(\alpha\), \[
q=M\sqrt{2\log(2T/\alpha)/n},\qquad
e_t=q[(1+\eta L)^t-1]/L.
\] With probability at least \(1-\alpha\), simultaneously through \(T\),
\[
|x_t-\bar x_t|\le e_t,\qquad
|r_p(x_t)-r_p(\bar x_t)|\le e_t/(\sigma\sqrt{2\pi}).
\] Conditional Hoeffding bounds followed by
\(d_{t+1}\le(1+\eta L)d_t+\eta q\) prove this. Its exponential
amplification can make it uninformative before rebound; it is a
baseline, not a principal result.

\subsection{8. Critical nonlinear SGD: a quantitative joint
distribution}\label{draft-critical-nonlinear-sgd-a-quantitative-joint-distribution}

Throughout this section fix \(0<m\le\sigma\), \(0<\eta<1/L\), and \[
0<\epsilon<1/2-\Phi(-m/\sigma).
\] Take \(|\rho|\le\rho_0\) for a sufficiently small fixed \(\rho_0\).
Constants may depend on these fixed parameters and on the initialization
bound, but not on \(n,t\) or the parameter sequence within this
neighborhood.

\subsubsection{Lemma 8. A normalized deterministic
conjugacy}\label{draft-lemma-8.-a-normalized-deterministic-conjugacy}

Let \(G_\rho(x)=x+\eta v_\rho(x)\) and
\(A_\rho=G_\rho'(q_\rho)=1+\eta\lambda_\rho>1\). There is a strictly
increasing \(C^2\) function \(H_\rho\) with \[
H_\rho(q_\rho)=0,\qquad H_\rho'(q_\rho)=1,\qquad
H_\rho(G_\rho(x))=A_\rho H_\rho(x).
\] On fixed compact intervals, \(H_\rho,H_\rho',H_\rho''\) are uniformly
bounded for small \(\rho\), \(H_\rho'\) has a uniform positive lower
bound, and these functions depend continuously on \(\rho\). At balance,
\(H_0\) is odd.

\textbf{Proof.} See \href{gaussian_proofs.md}{Gaussian proofs}, under
the same result number.

\subsubsection{Theorem 9. State-dependent SGD admits a Gaussian-seed
exit
approximation}\label{draft-theorem-9.-state-dependent-sgd-admits-a-gaussian-seed-exit-approximation}

Suppose \[
|x_0-q_\rho|\le K/\sqrt n
\] for a fixed \(K<\infty\). For sufficiently large \(n\), the safe set
is a finite interval \([a_n,c_n]\) containing \(q_\rho\) in its
interior, with peak-to-boundary distances uniformly positive. Put \[
\tau=T_\epsilon,\qquad
J=\begin{cases}+1,&x_\tau>c_n,\\-1,&x_\tau<a_n,\end{cases}
\] \[
\mu_n=\sqrt n H_\rho(x_0),\qquad
s_n^2=\frac{\eta^2\operatorname{Var}_\rho U(q_\rho,z)}{A_\rho^2-1},
\] \[
B_n^+=H_\rho(c_n)>0,\qquad B_n^-=-H_\rho(a_n)>0.
\] Then \(\tau<\infty\) almost surely. There is a constant \(C\) such
that, for every integer \(t\ge0\), \[
\left|\Pr(\tau\le t,J=+1)
-\Phi\!\left(\frac{\mu_n-\sqrt n A_\rho^{-t}B_n^+}{s_n}\right)\right|
\le C\left[n^{-1/2}+
\left(\frac{1+\log n}{n}+A_\rho^{-2t}\right)^{1/3}\right],
\] \[
\left|\Pr(\tau\le t,J=-1)
-\Phi\!\left(\frac{-\mu_n-\sqrt n A_\rho^{-t}B_n^-}{s_n}\right)\right|
\le C\left[n^{-1/2}+
\left(\frac{1+\log n}{n}+A_\rho^{-2t}\right)^{1/3}\right].
\] Also, \[
|\Pr(J=+1)-\Phi(\mu_n/s_n)|
\le C[n^{-1/2}+((1+\log n)/n)^{1/3}].
\] In the logarithmic window \(nA_\rho^{-2t}=O(1)\), the error is
\(O((\log n/n)^{1/3})\). This rate is not asserted optimal, and \(C\) is
not numerically calibrated.

\textbf{Proof.} See \href{gaussian_proofs.md}{Gaussian proofs}, under
the same result number.

\subsection{8A. General scalar laws and estimated
geometry}\label{draft-a.-general-scalar-laws-and-estimated-geometry}

The Gaussian model supplies explicit distributional geometry. The
extension in \href{scalar_exit.md}{general scalar exit theory} instead
fixes any score law \(Q\), including a finite weighted pool, and assumes
one repelling root on the required compact domain, positive root noise
variance, fixed positive step size, two finite interior exit boundaries
with an overshoot margin, and \(O(n^{-1/2})\) initialization relative to
the true root (G1--G4). These assumptions are not inferred from local
positive curvature alone. Lemmas E.1--E.2 and Theorem E.3 give the
actual entropy-SGD joint interval-exit approximation, including \[
\sup_{t\in\mathbb Z_{\ge0}}\max_\pm
|\Pr(\tau_n\le t,J_n=\pm1)-\Psi_n^\pm(t)|\le C e_n,
\qquad e_n=n^{-1/2}+((1+\log n)/n)^{1/3}.
\] The early-time stopped-martingale tail argument is essential for this
uniform statement. The original Theorem 9 retains its Gaussian-family
formulation and proof; the extension does not automatically imply
uniformity over changing score laws without common geometric constants.

\href{calibration.md}{Independent calibration theory} treats a separate
iid sample of size \(M\) from the same \(Q\). It changes the predictor,
not the actual initialization or adaptation recursion. Lemma E.4
controls normalized conjugacy under small \(C^3\) field perturbations;
Proposition E.5 gives empirical field concentration and the root CLT.
Theorem E.6 yields, on the declared good-calibration and endpoint event,
\[
\sup_t\max_\pm |\Pr(\tau_n\le t,J_n=\pm1)-\widehat\Psi_n^\pm(t)|
\le C[e_n+\sqrt n|d_q|+\sqrt n d_q^2
 +(1+\log n)d_v+d_\sigma+d_b],
\] where \(d_q=\widehat q-q\), \(d_v=\|\widehat v-v\|_{C^3}\),
\(d_\sigma=|\widehat\sigma_g^2-\sigma_g^2|\), and \(d_b\) is endpoint
error. Thus peak error acts on the shrinking SGD seed scale, growth-rate
error accumulates over logarithmic time, and endpoint error remains a
separate input. With fixed \(Q,w,\eta\), consistent endpoints and
\(M/n\to\infty\) suffice for consistency of this empirical-root
predictor; no universal sample-complexity lower bound is asserted.

For \(s^2=\eta^2\sigma_g^2/((1+\eta\lambda)^2-1)\), define \[
\mathcal R_{M,n}^2=\frac nM\frac{2+\eta\lambda}{\eta\lambda}.
\] If \(\sqrt n(x_{0,n}-q)/s\to z\) and
\(\mathcal R_{M,n}\to r<\infty\), Theorem E.7 gives \[
\Pr(J_n=+1)\to\Phi(z),\qquad
\widehat\pi_{M,n}\Rightarrow\Phi(z-rZ),\quad Z\sim N(0,1).
\] The randomness on the right is over calibration samples. At exact
critical centering, Corollary E.8 gives limiting mean absolute
direction-probability error \(\arctan(r)/\pi\), equal to \(1/4\) when
\(r=1\). Joint time limits retain an integer-phase subsequence. This
random forecast limit does not inflate the variance of the unchanged SGD
process. The distinct calibration-driven initialization result is
Proposition E.13 in \href{risk_boundaries.md}{the boundary and pool
part}.

\href{risk_boundaries.md}{Risk boundaries and empirical pools} supplies
the remaining interfaces. Proposition E.9 distinguishes a fixed off-peak
initialization, whose strict-margin exit converges to a finite
deterministic count, from a genuinely critical parameter sequence.
Proposition E.10 controls regular differentiable endpoints by
risk-difference error divided by a positive transversality constant; it
does not apply directly to step-function empirical risk. Proposition
E.11 adds a component-skip probability when transferring interval exit
to actual risk exit. Proposition E.12 separates conditional-pool
inference from population-risk uncertainty. None estimates
label-dependent risk endpoints from an arbitrary unlabeled score
marginal.

These are supplied statements with proofs, integrated here without
claiming a completed independent proof audit or numerical calibration of
their constants. The visual experiments must still verify interval-wide
conditions, initialization regime, endpoint conventions, and
component-skip discrepancies before any setting is declared covered. The
extension does not generalize the Gaussian inverse-batch-tolerance
result to arbitrary score laws.

\subsection{9. Critical scaling, lattice phase, and microscopic
tolerance}\label{draft-critical-scaling-lattice-phase-and-microscopic-tolerance}

\subsubsection{Corollary 10. The critical joint
limit}\label{draft-corollary-10.-the-critical-joint-limit}

Suppose \[
\sqrt n\rho_n\to c\in\mathbb R,\qquad
\sqrt n x_{0,n}\to u\in\mathbb R.
\] Define \[
\mu=u+\beta c/\lambda,\qquad
s_0^2=\operatorname{Var}_0U(0,z)>0,
\] \[
s_\eta^2=\frac{\eta^2s_0^2}{A^2-1}
=\frac{\eta s_0^2}{\lambda(2+\eta\lambda)}.
\] Let \(c_\epsilon>0\) solve \(r_0(c_\epsilon)=r_0(0)+\epsilon\), and
set \(C_\epsilon=H_0(c_\epsilon)>0\). Then \[
\Pr(J=+1)\to\Phi(\mu/s_\eta).
\] For integer time sequences with
\(\sqrt n A^{-t_n}\to r\in(0,\infty)\), \[
\Pr(T_\epsilon\le t_n,J=+1)\to
\Phi((\mu-C_\epsilon r)/s_\eta),
\] \[
\Pr(T_\epsilon\le t_n,J=-1)\to
\Phi((-\mu-C_\epsilon r)/s_\eta),
\] and \[
T_\epsilon=\frac{\log n}{2\log A}+O_{\mathbb P}(1).
\]

\textbf{Proof.} See \href{gaussian_proofs.md}{Gaussian proofs}, under
the same result number.

The lattice cannot generally be removed at fixed step size. Put \[
b_n=\log\sqrt n/\log A.
\] Along a subsequence whose fractional part satisfies
\(\{b_n\}\to\theta\), for integer \(k\), \[
\Pr(T_\epsilon-\lfloor b_n\rfloor\le k)
\to\Pr(|N(\mu,s_\eta^2)|>C_\epsilon A^{\theta-k}).
\] Equivalently, the integer-valued limit is \[
\left\lceil\theta+
\frac{\log C_\epsilon-\log|N(\mu,s_\eta^2)|}{\log A}\right\rceil.
\] Gaussian boundary atoms have probability zero. The time fluctuation
is a rounded negative logarithm of a Gaussian magnitude; it is not
Gaussian and is not automatically Gumbel.

The effective seed signal-to-noise ratio is \[
\kappa_n\simeq\frac{\sqrt n(x_0-q_\rho)}{s_\eta}
\simeq\frac{\sqrt n(\beta\rho/\lambda-\delta)}{s_\eta}.
\] Translation can cancel or reinforce the leading prior seed. For pure
prior shift, \[
\rho_{\rm cross}\simeq\frac{\lambda s_\eta}{\beta\sqrt n}.
\] If \(\rho_n\downarrow0\), \(\sqrt n\rho_n\to\infty\),
\(x_0=\alpha_0\rho_n\), and \(\alpha_0+\beta/\lambda\ne0\), the seed
dominates noise and \[
T_\epsilon=\log(1/\rho_n)/\log A+O_{\mathbb P}(1).
\] To see this, use the same weighted conjugacy expansion normalized by
\(\rho_n\): the noise is
\(O_{\mathbb P}(n^{-1/2})=o_{\mathbb P}(\rho_n)\), the Taylor remainder
is \(O_{\mathbb P}(n^{-1})\), and the initial seed is
\((\alpha_0+\beta/\lambda)\rho_n+o(\rho_n)\). The sign is consequently
deterministic in the limit. Leading-seed cancellation requires another
scale analysis.

Although \(s_\eta\sim s_0\sqrt{\eta/(2\lambda)}\) algebraically as
\(\eta\downarrow0\), all probability limits above fix \(\eta>0\). They
do not prove a joint limit with \(\eta=\eta_n\to0\). Likewise the
logarithmic batch-size delay applies to balanced or critically rescaled
families, not to every fixed nonzero prior.

\subsubsection{Theorem 11. Inverse-batch tolerance gives finite-step AR
exit}\label{draft-theorem-11.-inverse-batch-tolerance-gives-finite-step-ar-exit}

Retain the critical scaling and let \(n\epsilon_n\to e\in[0,\infty)\).
Define \[
X_0=u,\qquad X_{t+1}=AX_t+\eta\beta c+\eta s_0Z_t,
\qquad Z_t\overset{\rm iid}{\sim}N(0,1),
\] and \[
I_{c,u,e}=
[k_0c-\sqrt{(u-k_0c)^2+2e/B},\
 k_0c+\sqrt{(u-k_0c)^2+2e/B}],
\] where \(B=m\varphi(m/\sigma)/\sigma^3\). For
\(\tau_{\rm AR}=\inf\{t\ge1:X_t\notin I_{c,u,e}\}\), \[
T_{\epsilon_n}\Rightarrow\tau_{\rm AR},\qquad
T_{\epsilon_n}=O_{\mathbb P}(1).
\]

\textbf{Proof.} See \href{gaussian_proofs.md}{Gaussian proofs}, under
the same result number.

\subsubsection{Proposition 12. Zero tolerance and a limitation of
shrinking the
step}\label{draft-proposition-12.-zero-tolerance-and-a-limitation-of-shrinking-the-step}

If \(\rho=0,x_0=0,\epsilon=0\), then for every finite \(n\) and every
\(\eta>0\), \[
T_0=1\quad\text{almost surely}.
\] For pure prior shift \(x_0=0,\sqrt n\rho_n\to c>0\), every
deterministic sequence \(\eta_n>0\) satisfies \[
\liminf_n\Pr(r_{\rho_n}(x_1)>r_{\rho_n}(0))
\ge\Phi(-\beta c/s_0)>0.
\] For fixed \(\eta>0\), the first-step improvement probability
converges to \[
\Phi\!\left(\frac{(2k_0-\eta\beta)c}{\eta s_0}\right)
-\Phi(-\beta c/s_0).
\]

\textbf{Proof.} See \href{gaussian_proofs.md}{Gaussian proofs}, under
the same result number.

The limitation concerns unconditional execution of the empirical entropy
direction. It does not exclude rejecting an update, using extra
information, or a different unlabeled adaptation method. The probability
of a strict risk increase also says nothing by itself about the size of
that increase.

\subsection{10. Numerical evidence}\label{draft-numerical-evidence}

The experimental protocols, tables, ablations, censoring conventions,
and numerical limitations are in \href{experiments.md}{the experimental
part}. Sections 10.6--10.9 distinguish mixed-shift evidence, failed
source-calibrated visual prediction, target empirical geometry
diagnostics, and independent unlabeled field calibration.

\subsection{11. What would make a stopping rule
observable?}\label{draft-what-would-make-a-stopping-rule-observable}

The risk boundaries used above depend on the target parameters and serve
as analytical objects. They do not themselves define an implementable
unlabeled rule.

A generic certificate clarifies the additional requirement. Let
\(\vartheta=(\delta,p)\), with known anchor parameters, and suppose
label-free, nested confidence sets \(\mathcal C_t\) satisfy \[
\Pr(\vartheta\in\mathcal C_t\text{ for all }t)\ge1-\alpha.
\] For a proposed threshold \(b'\), accept only if \[
\sup_{\vartheta\in\mathcal C_t}
\{R_\vartheta(b')-R_\vartheta(b_0)\}\le\epsilon.
\] On the coverage event every accepted iterate is safe, including
data-dependent proposals. This follows directly from the supremum and is
a generic baseline.

To become a useful algorithm, this construction needs a sequential
confidence-set method under the stated observation model, accounting for
repeated inspection and proposal dependence. It also needs a nontrivial
acceptance and improvement guarantee. Ordinary fixed-time intervals
repeatedly reused do not suffice. Anchor error, unknown component
scales, orientation mistakes, and model misspecification require
explicit treatment.

Appropriate empirical comparators include frozen source, original
entropy updates, logit-adjusted entropy updates, mean recentering, an
oracle risk threshold, and one-time unlabeled calibration with a frozen
plug-in estimate. Known anchors supply more information than strict
weights-only source-free deployment. No universal claim about all
unlabeled adaptation is implied.

\subsection{12. Limitations and falsifiable
scope}\label{draft-limitations-and-falsifiable-scope}

The integrated extension supplies proofs for independent-geometry
prediction error and regular boundary perturbations under explicit
assumptions. Independent verification and alignment with visual events
remain necessary. The expanded calibration grid now covers 144 cells
with 20 repetitions, including nonbalanced priors and nonzero offsets,
but omits offset +2 and does not exactly test target-centered critical
limits. Noise ablations support continuing far-field noise as an
important source of tail dispersion in three single-side-dominated
settings; a combined early-seed and passage-noise prediction for mixed
directions remains open. Unlabeled risk-boundary identification, the
exact separation phase boundary, numerical approximation constants, and
unknown-anchor robustness remain unresolved.

The completed frozen-encoder study retains official train/test splits
and controls projected translation and priors. It supplies useful
conditional-oracle evidence, but the original independently calibrated
prediction fails overall and the improved field-calibration study still
uses labeled oracle risk endpoints for its strongest results. Neither
full BN adaptation nor scale/LA comparisons have yet been validated by
the available results. Image-space corruptions should not be claimed to
satisfy a Gaussian projection model without checking that approximation.
The general scalar theorem now states such geometric conditions; their
validity on each visual exit domain remains to be checked.

The model omits nonstationarity, batch reuse, temporal dependence,
learning-rate schedules, momentum, joint feature learning, and
multilayer normalization dynamics. A safe population basin is not
permanent finite-batch safety. A high right-exit probability is not a
safety probability. A fixed-tolerance logarithmic horizon is not an
improvement guarantee, and a longer horizon under larger batch size is
not universal at fixed prior.

The contribution would be insufficient if its statements followed
immediately from accepted entropy-drift results plus a generic
unstable-exit theorem without a meaningful new risk prediction. The
standard mechanisms are acknowledged explicitly. The case for a
contribution must instead rest on the actual recursion, risk-defined
boundary, finite-parameter predictions, and tolerance distinction,
followed by a direct comparison with existing results.

\subsection{References}\label{draft-references}

{[}1{]} T. Weerasinghe, R. Tennakoon, W. Chuah, and A. Bab-Hadiashar.
\emph{Logit-Adjusted Test-Time Adaptation under Partial Class
Imbalance}. WACV, 2026.
\href{https://openaccess.thecvf.com/content/WACV2026/papers/Weerasinghe_Logit-Adjusted_Test-Time_Adaptation_under_Partial_Class_Imbalance_WACV_2026_paper.pdf}{Official
paper}.

{[}2{]} Hoang, Vo, and Do. \emph{Persistent Test-time Adaptation in
Recurring Testing Scenarios}. arXiv:2311.18193, v4, 2024.
\href{https://arxiv.org/abs/2311.18193}{Paper}.

{[}3{]} Sahoo et al.~\emph{A Layer Selection Approach to Test Time
Adaptation}. arXiv:2404.03784, v2, 2025. Earlier version:
\emph{Layerwise Early Stopping for Test Time Adaptation}.
\href{https://arxiv.org/abs/2404.03784}{Paper}.

{[}4{]} Shin et al.~\emph{STAG}. arXiv:2402.09004, v2, 2026.
\href{https://arxiv.org/abs/2402.09004}{Paper identifier}.

{[}5{]} Zhou et al.~\emph{On the Learnability of Test-Time Adaptation: A
Recovery Complexity Perspective}. arXiv:2605.28057, v2, 2026.
\href{https://arxiv.org/abs/2605.28057}{Paper}.

{[}6{]} Jhawar and Wang. \emph{When is Test-Time Adaptation Identifiable
From Unlabeled Evidence?} arXiv:2609.11235, v1, 2026.
\href{https://arxiv.org/abs/2609.11235}{Paper}.

{[}7{]} A. Saumard and J. A. Wellner. \emph{Log-concavity and strong
log-concavity: a review}. Statistics Surveys 8, 2014.
\href{https://arxiv.org/abs/1404.5886}{Paper}.

{[}8{]} Y. Kifer. \emph{The exit problem for small random perturbations
of dynamical systems with a hyperbolic fixed point}. Israel Journal of
Mathematics 40:74--96, 1981.
\href{https://cris.huji.ac.il/en/publications/the-exit-problem-for-small-random-perturbations-of-dynamical-syst/}{Publication
record}.

{[}9{]} Y. Bakhtin. \emph{Exit asymptotics for small diffusion about an
unstable equilibrium}. arXiv:math/0701569, 2007. DOI:
10.1016/j.spa.2007.06.003.
\href{https://arxiv.org/abs/math/0701569}{Paper}.

{[}10{]} Y. Bakhtin and A. Gaál. \emph{Exit time asymptotics for
dynamical systems with fast random switching near an unstable
equilibrium}. arXiv:1901.05513, 2019.
\href{https://arxiv.org/abs/1901.05513}{Paper}.

{[}11{]} Yarin Bar, Shalev Shaer and Yaniv Romano. \emph{Protected
Test-Time Adaptation via Online Entropy Matching: A Betting Approach}.
NeurIPS 2024.
\href{https://proceedings.neurips.cc/paper_files/paper/2024/file/9b35a0a20d617dc68ae98a7a57df2f51-Paper-Conference.pdf}{Official
paper}.
